\chapter{Kohomologia singularna i iloczyn kubkowy}
\label{ch:kohomologia-singularna}

Homologia przypisuje przestrzeni kompleks łańcuchowy i bierze jądra
modulo obrazy różniczek. W kohomologii stosujemy do łańcuchów
kontrawariantny funktor homomorfizmów do grupy współczynników.
W rezultacie uzyskujemy nie tylko grupy, lecz także naturalne mnożenie.
Ten rozdział łączy konstrukcje homologiczne z dualnością Poincarégo
i porównaniem de Rhama omawianym dalej.

\section{Kołańcuchy, koróżniczka i funktorialność}

\begin{definicja}
Niech $X$ będzie przestrzenią topologiczną, $G$ grupą abelową i
$(C_n(X),\partial_n)$ singularnym kompleksem łańcuchowym z całkowitymi
współczynnikami. Dla $n\geq0$ określamy grupę \emph{kołańcuchów}
$C^n(X;G)=\operatorname{Hom}_{\Z}(C_n(X),G)$, a dla $n<0$ przyjmujemy
$C^n(X;G)=0$. Homomorfizm
$\delta^n:C^n(X;G)\to C^{n+1}(X;G)$ jest dany wzorem
$(\delta^n\varphi)(c)=\varphi(\partial_{n+1}c)$.
Grupę kohomologii stopnia $n$ definiujemy przez
\[
H^n(X;G)=\ker\delta^n/\operatorname{im}\delta^{n-1}.
\]
Elementy $\ker\delta^n$ są kocyklami, a elementy
$\operatorname{im}\delta^{n-1}$ kobrzegami.
\end{definicja}

\begin{stwierdzenie}
Dla każdej przestrzeni topologicznej $X$ i grupy abelowej $G$
zachodzi $\delta^{n+1}\circ\delta^n=0$ dla wszystkich $n\geq0$.
\end{stwierdzenie}
\begin{proof}
Dla $\varphi\in C^n(X;G)$ i $c\in C_{n+2}(X)$ mamy
$(\delta^{n+1}\delta^n\varphi)(c)
=\varphi(\partial_{n+1}\partial_{n+2}c)=0$,
ponieważ $\partial^2=0$ w kompleksie łańcuchów singularnych.
Zatem każdy kobrzeg jest kocyklem i iloraz w definicji ma sens.
\end{proof}

\begin{definicja}
Niech $f:X\to Y$ będzie ciągłe. Dla $\varphi\in C^n(Y;G)$
określamy $(f^*\varphi)(c)=\varphi(f_\#c)$ dla $c\in C_n(X)$,
gdzie $f_\#$ jest odwzorowaniem łańcuchowym indukowanym przez $f$.
\end{definicja}

\begin{twierdzenie}[Kontrawariantność i niezmienniczość homotopijna]
Każde ciągłe $f:X\to Y$ indukuje homomorfizmy
$f^*:H^n(Y;G)\to H^n(X;G)$, dla których
$(g\circ f)^*=f^*\circ g^*$, a $(\id_X)^*=\id$.
Jeżeli $f\simeq g$, to $f^*=g^*$.
\end{twierdzenie}
\begin{proof}
Tożsamość $\partial f_\#=f_\#\partial$ po przejściu do
$\operatorname{Hom}(-,G)$ daje $\delta f^*=f^*\delta$.
Stąd $f^*$ odwzorowuje kocykle w kocykle i kobrzegi w kobrzegi.
Tożsamości związane ze składaniem wynikają już na poziomie łańcuchów.
Jeśli $f\simeq g$, w dowodzie niezmienniczości homotopijnej
homologii skonstruowano operator pryzmowy $P$ spełniający
$g_\#-f_\#=\partial P+P\partial$.
Po złożeniu z kołańcuchem $\varphi$ otrzymujemy
$g^*\varphi-f^*\varphi=\delta(P^*\varphi)+P^*(\delta\varphi)$,
gdzie $P^*$ obniża stopień o jeden. Dla kocyklu drugi składnik znika,
a pierwszy jest kobrzegiem. Indukowane odwzorowania są równe.
\end{proof}

\begin{przyklad}
Dla niepustej łukowo spójnej przestrzeni $X$ mamy $H^0(X;G)\cong G$:
kocykl stopnia zero jest funkcją na punktach, która przyjmuje tę samą
wartość na końcach każdej drogi. Nie ma nietrywialnych kobrzegów
stopnia zero. Dla przestrzeni o $k$ składowych łukowych otrzymujemy
$H^0(X;G)\cong G^k$; przy nieskończenie wielu składowych
łukowych pojawia się iloczyn, nie suma prosta.
\end{przyklad}

\section{Kohomologia względna i ciąg dokładny}

\begin{definicja}
Dla pary przestrzeni $A\subseteq X$ przyjmujemy
$C^n(X,A;G)=\{\varphi\in C^n(X;G):\varphi|_{C_n(A)}=0\}$.
Ponieważ brzeg łańcucha z $A$ pozostaje w $A$, jest to podkompleks
kołańcuchowy. Definiujemy $H^n(X,A;G)=H^n(C^*(X,A;G))$.
\end{definicja}

\begin{twierdzenie}[Długi ciąg dokładny pary]
Dla każdej pary przestrzeni topologicznych $A\subseteq X$ oraz
każdej grupy abelowej $G$ istnieje naturalny długi ciąg dokładny
\[
\cdots\to H^n(X,A;G)\to H^n(X;G)\to H^n(A;G)
\xrightarrow{\,\partial^*\,}H^{n+1}(X,A;G)\to\cdots.
\]
\end{twierdzenie}
\begin{proof}
Łańcuchy singularne są wolnymi grupami abelowymi na singularnych
sympleksach. Każdy sympleks w $A$ jest też sympleksem w $X$, więc
$C_n(A)$ jest wolnym składnikiem prostym $C_n(X)$.
Wobec tego ograniczenie kołańcuchów z $X$ do $A$ jest surjekcją
$C^n(X;G)\to C^n(A;G)$, a jej jądrem jest $C^n(X,A;G)$.
Otrzymujemy krótki ciąg dokładny kompleksów kołańcuchowych.
Konstrukcja homomorfizmu łączącego, przedstawiona w rozdziale
o algebrze homologicznej, daje wskazany długi ciąg. Naturalność
wynika z naturalności restrykcji i konstrukcji homomorfizmu łączącego.
\end{proof}

\section{Iloczyn kubkowy: skąd bierze się mnożenie klas}

\begin{definicja}[Iloczyn kubkowy Alexandra--Whitneya]
Niech $R$ będzie przemiennym pierścieniem z jedynką,
$\alpha\in C^p(X;R)$ i $\beta\in C^q(X;R)$, gdzie $p,q\geq0$.
Dla singularnego sympleksu $\sigma:\Delta^{p+q}\to X$ definiujemy
\[
(\alpha\smile\beta)(\sigma)=
\alpha\bigl(\sigma|_{[v_0,\ldots,v_p]}\bigr)\,
\beta\bigl(\sigma|_{[v_p,\ldots,v_{p+q}]}\bigr).
\]
Restrykcje do ścian są rozumiane przez standardowe liniowe
parametryzacje sympleksów o wskazanych wierzchołkach.
Definicję przedłużamy liniowo na wszystkie łańcuchy.
\end{definicja}

\begin{stwierdzenie}[Reguła Leibniza]
Jeżeli $\alpha\in C^p(X;R)$ i $\beta\in C^q(X;R)$, to
$\delta(\alpha\smile\beta)
=(\delta\alpha)\smile\beta+(-1)^p\alpha\smile(\delta\beta)$.
W konciągu iloczyn kocykli jest kocyklem, a gdy co najmniej
jeden czynnik jest kobrzegiem, iloczyn jest kobrzegiem.
\end{stwierdzenie}
\begin{proof}
Obliczamy lewą stronę na $\sigma:\Delta^{p+q+1}\to X$,
rozwijając $\partial\sigma$ jako naprzemienną sumę ścian.
Oznaczmy $N=p+q+1$ i dla czytelności zapisu pomijajmy
symbol $\sigma$ w restrykcjach do ścian.
Po obliczeniu lewej strony z definicji różniczki dostajemy
\[
\begin{aligned}
\delta(\alpha\smile\beta)(\sigma)
={}&\sum_{i=0}^{p}(-1)^i
 \alpha([v_0,\ldots,\widehat v_i,\ldots,v_{p+1}])
 \beta([v_{p+1},\ldots,v_N])\\
&+\sum_{i=p+1}^{N}(-1)^i
 \alpha([v_0,\ldots,v_p])
 \beta([v_p,\ldots,\widehat v_i,\ldots,v_N]).
\end{aligned}
\]
Istotnie, dla $i\leq p$ po usunięciu jednego
wierzchołka pierwszy odcinek ma nadal $p+1$ wierzchołków;
dla $i\geq p+1$ pierwszy odcinek pozostaje bez zmian.
Z drugiej strony $(\delta\alpha)\smile\beta$
zawiera \emph{pierwszą} sumę, a dodatkowo jej wyraz
o indeksie $i=p+1$, równy
$(-1)^{p+1}\alpha([v_0,\ldots,v_p])
\beta([v_{p+1},\ldots,v_N])$.
Podobnie $(-1)^p\alpha\smile(\delta\beta)$
zawiera \emph{drugą} sumę, a dodatkowo wyraz
o indeksie $i=p$ (tj. indeksie $0$ w drugiej ścianie),
równy
$(-1)^p\alpha([v_0,\ldots,v_p])
\beta([v_{p+1},\ldots,v_N])$.
Te dwa dodatkowe wyrazy znoszą się.
Otrzymujemy żądaną tożsamość Leibniza.
Dla $p=q=0$ sprowadza się ona do zwykłej równości
$\alpha(v_1)\beta(v_1)-\alpha(v_0)\beta(v_0)
=(\alpha(v_1)-\alpha(v_0))\beta(v_1)
+\alpha(v_0)(\beta(v_1)-\beta(v_0))$.
Jeżeli $\delta\alpha=0$ i $\beta=\delta\gamma$, to
$\alpha\smile\beta=(-1)^p\delta(\alpha\smile\gamma)$;
przypadek kobrzegu w pierwszym czynniku jest analogiczny.
\end{proof}

\begin{twierdzenie}
Dla każdej przestrzeni $X$ i przemiennego pierścienia $R$ iloczyn
kubkowy indukuje dwuliniowe odwzorowania
$H^p(X;R)\times H^q(X;R)\to H^{p+q}(X;R)$.
Z tym mnożeniem $H^*(X;R)=\bigoplus_{n\geq0}H^n(X;R)$ jest
łącznym, unitalnym, \emph{stopniowo przemiennym} pierścieniem,
tzn. dla jednorodnych $a,b$ zachodzi
$a\smile b=(-1)^{\deg(a)\deg(b)}b\smile a$.
Ponadto $f^*(a\smile b)=f^*a\smile f^*b$ dla ciągłego $f$.
\end{twierdzenie}
\begin{proof}[Argument z użyciem standardowej homotopii diagonalnych]
Poprawność definicji na kohomologii wynika z reguły Leibniza.
Łączność wynika bezpośrednio z trzykrotnego rozcięcia
uporządkowanego zbioru wierzchołków sympleksu.
Elementem jednostkowym jest klasa stałego kocyklu
$1\in C^0(X;R)$, przyjmującego wartość $1_R$ na każdym
wierzchołku. Naturalność wynika z faktu, że złożenie z $f$
komutuje z ograniczeniem sympleksu do ściany.
Stopniowa przemienność nie zachodzi na ogół na poziomie
kołańcuchów: jej dowód wymaga homotopii między
odwzorowaniem diagonalnym Alexandra--Whitneya a jego
transpozycją. Istnienie takiej homotopii wynika z kołańcuchowego lematu
o porównaniu diagonali (jego niezależny dowód wymaga jawnej
konstrukcji homotopii pryzmowej). Daje ona
$\alpha\smile\beta-(-1)^{pq}\beta\smile\alpha$ jako kobrzeg,
gdy $\alpha$ i $\beta$ są kocyklami.
\end{proof}

\begin{przyklad}[Torus: jawny rachunek iloczynu kubkowego]
\label{prz:kohomologia-kubek-torus}
Niech $a,b\in H^1(T^2;\Z)$ będą cofnięciami generatora
$H^1(S^1;\Z)\cong\Z$ przez obie projekcje
$T^2\to S^1$. Wówczas $a\smile a=b\smile b=0$:
każdy z tych iloczynów jest cofnięciem klasy
w $H^2(S^1;\Z)=0$. Wybierzmy natomiast
standardowy kwadrat fundamentalny torusa
o wierzchołkach
$A=(0,0)$, $B=(1,0)$, $C=(1,1)$ i $D=(0,1)$,
i podzielmy go przekątną $AC$ na dwa
zorientowane trójkąty $[A,B,C]$ oraz $[A,C,D]$.
Ich suma reprezentuje klasę fundamentalną torusa:
wewnętrzne krawędzie $AC$ kasują się,
a przeciwległe zewnętrzne krawędzie
są identyfikowane z przeciwnymi orientacjami
w ich obiegach brzegowych.

Przedstawmy $a,b$ przez kocykle rejestrujące
odpowiednio poziome i pionowe przejścia
w nakryciu $\R^2\to T^2$. Na krawędziach
tej triangulacji ich wartości są następujące:
\[
\begin{array}{c|rrrrr}
 & AB & BC & CD & AD & AC\\ \hline
 a & 1&0&-1&0&1\\
 b & 0&1&0&1&1
\end{array}.
\]
Wartości na przekątnej wynikają z połączenia
jednego przesunięcia poziomego i jednego pionowego;
równania kocykli na obu trójkątach można sprawdzić
bezpośrednio, np. $a(BC)-a(AC)+a(AB)=0-1+1=0$.
Wzór Alexandra--Whitneya daje
\[
 (a\smile b)([A,B,C])=a(AB)b(BC)=1,\qquad
 (a\smile b)([A,C,D])=a(AC)b(CD)=0.
\]
Wartość iloczynu na klasie fundamentalnej wynosi więc $1$.
Oznacza to, że $a\smile b$ jest generatorem
$H^2(T^2;\Z)\cong\Z$.
Stopniowa przemienność daje
$b\smile a=-a\smile b$.
Ostatecznie pierścień kohomologii torusa
jest algebrą zewnętrzną na generatorach $a,b$ stopnia $1$.
Sama znajomość grup $H^0\cong\Z$,
$H^1\cong\Z^2$, $H^2\cong\Z$
nie ujawnia jeszcze tej struktury mnożenia.
\end{przyklad}

\begin{figure}[!htbp]
\centering
\begin{tikzpicture}[>=Stealth,font=\small]
  \fill[manifoldblue!7] (0,0) rectangle (4.6,3.2);
  \fill[manifoldgreen!15] (0,0)--(4.6,0)--(4.6,3.2)--cycle;
  \fill[manifoldblue!14] (0,0)--(4.6,3.2)--(0,3.2)--cycle;
  \draw[thick,manifoldblue!75!black] (0,0) rectangle (4.6,3.2);
  \draw[thick] (0,0)--(4.6,3.2);
  \draw[->,manifoldred,very thick] (.55,-.10)--(3.9,-.10);
  \draw[->,manifoldred,very thick] (.55,3.3)--(3.9,3.3);
  \draw[->,manifoldgold!75!black,very thick] (-.10,.55)--(-.10,2.75);
  \draw[->,manifoldgold!75!black,very thick] (4.7,.55)--(4.7,2.75);
  \node[below left] at (0,0) {$A$};
  \node[below right] at (4.6,0) {$B$};
  \node[above right] at (4.6,3.2) {$C$};
  \node[above left] at (0,3.2) {$D$};
  \node at (3.35,.95) {$[A,B,C]$};
  \node at (1.40,2.35) {$[A,C,D]$};
  \node[manifoldred,below=5pt] at (2.3,0) {$a$};
  \node[manifoldgold!75!black,right=4pt] at (4.6,1.6) {$b$};
\end{tikzpicture}
\caption{Triangulacja kwadratu fundamentalnego torusa.
Iloczyn kubkowy $a\smile b$ przyjmuje na dwóch
zorientowanych trójkątach wartości odpowiednio $1$ i $0$.}
\label{fig:kubek-torus}
\end{figure}

\section{Współczynniki i twierdzenie o współczynnikach uniwersalnych}

\begin{uwaga}[Co oznacza $\operatorname{Ext}^1$?]
Dla grupy abelowej $A$ wybieramy dokładny ciąg
$0\to F_1\xrightarrow{u} F_0\to A\to0$ z wolnymi grupami
abelowymi $F_0,F_1$ i definiujemy
$\operatorname{Ext}^1_{\Z}(A,G)=
\operatorname{coker}\bigl(\operatorname{Hom}(F_0,G)
\xrightarrow{u^*}\operatorname{Hom}(F_1,G)\bigr)$.
Wynik nie zależy, z dokładnością do kanonicznego izomorfizmu,
od wyboru wolnej prezentacji. Na przykład
$\operatorname{Ext}^1_{\Z}(\Z/m\Z,G)\cong G/mG$,
co wynika z prezentacji $0\to\Z\xrightarrow{\cdot m}\Z
\to\Z/m\Z\to0$.
\end{uwaga}

\begin{twierdzenie}[Współczynniki uniwersalne dla kohomologii]
Niech $X$ będzie przestrzenią topologiczną, $G$ grupą abelową,
a $n\geq0$ liczbą całkowitą. Istnieje naturalny krótki ciąg
dokładny grup abelowych
\[
0\longrightarrow\operatorname{Ext}^1_{\Z}(H_{n-1}(X;\Z),G)
\longrightarrow H^n(X;G)
\longrightarrow\operatorname{Hom}_{\Z}(H_n(X;\Z),G)
\longrightarrow0,
\]
przy czym $H_{-1}(X;\Z)=0$. Ciąg rozszczepia się jako ciąg
grup abelowych, lecz nie ma na ogół naturalnego wyboru rozszczepienia.
\end{twierdzenie}



\begin{proof}
Piszemy $C_n=C_n(X;\Z)$,
$Z_n=\ker\partial_n$, $B_n=\operatorname{im}\partial_{n+1}$ oraz
$H_n=Z_n/B_n$. Wszystkie $C_n,Z_n,B_n$ są wolnymi grupami
abelowymi: $C_n$ z definicji, a pozostałe jako podgrupy
grup wolnych. Ponieważ $B_{n-1}$ jest wolna, krótki ciąg
$0\to Z_n\to C_n\xrightarrow{\partial_n}B_{n-1}\to0$
rozszczepia się. Można więc wybrać
$C_n\cong Z_n\oplus B_{n-1}$, a w tych współrzędnych
różniczka jest projekcją na drugi składnik, po której
następuje włączenie do $C_{n-1}$.

Dla kocyklu $\varphi\in\operatorname{Hom}(C_n,G)$ restrykcja
$\varphi|_{Z_n}$ znika na $B_n$, a więc definiuje
homomorfizm $H_n\to G$. Powstaje naturalna mapa
$\rho:H^n(X;G)\to\operatorname{Hom}(H_n,G)$,
$\rho[\varphi]([z])=\varphi(z)$. Jest dobrze określona:
kobrzeg znika na cyklach, a kocykl na brzegach.
Mapa $\rho$ jest surjektywna: dowolny homomorfizm
$H_n\to G$ cofamy do $Z_n$ i przedłużamy do
$C_n=Z_n\oplus B_{n-1}$ przez zero na drugim składniku.

Przedstawiciel elementu $\ker\rho$ może być wybrany
jako kocykl znikający na $Z_n$. Rzeczywiście, warunek
$\rho[\varphi]=0$ oznacza dokładnie $\varphi|_{Z_n}=0$.
Takie kocykle są równoważne homomorfizmom
$B_{n-1}=C_n/Z_n\to G$. Kobrzegi odpowiadają tym
homomorfizmom $B_{n-1}\to G$, które są restrykcjami
homomorfizmów $C_{n-1}\to G$. Ponieważ
$C_{n-1}\cong Z_{n-1}\oplus B_{n-2}$, otrzymujemy
\[
\ker\rho\cong
\operatorname{coker}\!\left(
\operatorname{Hom}(Z_{n-1},G)\longrightarrow
\operatorname{Hom}(B_{n-1},G)\right).
\]
Dokładny ciąg
$0\to B_{n-1}\to Z_{n-1}\to H_{n-1}\to0$
jest wolną prezentacją $H_{n-1}$, więc ostatni
kokernel to $\operatorname{Ext}^1_{\Z}(H_{n-1},G)$.
Dostajemy żądany krótki ciąg dokładny.

Wybrane rozszczepienia $C_n=Z_n\oplus B_{n-1}$
rozkładają grupę kocykli na część odpowiadającą
$\operatorname{Hom}(H_n,G)$ i część odpowiadającą
$\operatorname{Hom}(B_{n-1},G)$; kobrzegi leżą tylko
w drugiej. To zapewnia rozszczepienie otrzymanego
ciągu, ale wybór dopełnień $Z_n$ w $C_n$ nie jest
naturalny względem odwzorowań łańcuchowych.
\end{proof}

\begin{przyklad}[Torsja widoczna dopiero w kohomologii]
Dla $\R P^2$ mamy $H_1(\R P^2;\Z)\cong\Z/2\Z$
i $H_2(\R P^2;\Z)=0$. Z twierdzenia wynika
$H^1(\R P^2;\Z)=0$ oraz
$H^2(\R P^2;\Z)\cong\Z/2\Z$.
Nie wolno więc utożsamiać kohomologii integralnej
z dualem homologii integralnej.
\end{przyklad}

\section{Iloczyny przestrzeni: Eilenberg--Zilber i Künneth}
\label{sec:ez-kunneth-v4}

Samo znajdowanie grup $H_p(X)$ i $H_q(Y)$ nie wystarcza do
obliczenia $H_n(X\times Y)$: przy współczynnikach całkowitych
pojawia się poprawka torsyjna. Zaczniemy więc od \emph{kompleksów},
a nie od grup homologii.

\begin{definicja}[Iloczyn tensorowy kompleksów i funktor $\operatorname{Tor}$]
Dla nieujemnie stopniowanych kompleksów $C,D$ nad $\Z$ niech
$(C\otimes D)_n=\bigoplus_{p+q=n}C_p\otimes_\Z D_q$ i
$\partial(c\otimes d)=\partial c\otimes d+(-1)^p c\otimes\partial d$
przy $c\in C_p$. Bezpośredni rachunek daje $\partial^2=0$.
Jeśli $A,B$ są grupami abelowymi i
$0\to F_1\xrightarrow{u}F_0\to A\to0$ jest wolną prezentacją,
to $\operatorname{Tor}^{\Z}_1(A,B)$ oznacza
$\ker(u\otimes\id_B)$. Niezależność od prezentacji wynika
z porównywania jej z innymi prezentacjami poprzez podniesienia
homomorfizmów wolnych grup.
\end{definicja}

\begin{twierdzenie}[Eilenberg--Zilber, łącznie z homotopiami]
\label{thm:ez-v4}
Dla dowolnych przestrzeni topologicznych $X,Y$ istnieją naturalne
odwzorowania łańcuchowe
$\mathrm{AW}:C_*(X\times Y;\Z)\to C_*(X;\Z)\otimes C_*(Y;\Z)$
oraz $\mathrm{sh}$ w przeciwną stronę, takie że
$\mathrm{sh}\,\mathrm{AW}\simeq\id$ i
$\mathrm{AW}\,\mathrm{sh}\simeq\id$ jako mapy łańcuchowe.
\end{twierdzenie}
\begin{proof}
\emph{Krok 1: odwzorowanie Alexandra--Whitneya.}
Dla $\sigma=(\sigma_X,\sigma_Y):\Delta^n\to X\times Y$
przyjmujemy
\[
\mathrm{AW}(\sigma)=\sum_{i=0}^n
\sigma_X|[0,\ldots,i]\otimes\sigma_Y|[i,\ldots,n].
\]
Ograniczenia mają standardową rosnącą kolejność wierzchołków.
Sprawdźmy różniczkę. W $\partial\mathrm{AW}(\sigma)$
usunięcie wierzchołka przed indeksem rozcięcia odpowiada
usunięciu tej samej ściany z $\mathrm{AW}(\partial\sigma)$;
podobnie po rozcięciu, gdzie znak tensoru daje $(-1)^i$.
Wyrazy związane z usunięciem samego punktu rozcięcia
występują dwa razy z przeciwnymi znakami. Pozostają zatem
dokładnie wyrazy $\mathrm{AW}(\partial\sigma)$.

\emph{Krok 2: przeplot.}
Niech $\sigma:\Delta^p\to X$, $\tau:\Delta^q\to Y$.
$(p,q)$-przeplot to słowo długości $p+q$ o $p$ krokach
poziomych i $q$ pionowych, dające wierzchołki drogi kratowej
$(a_j,b_j)$ od $(0,0)$ do $(p,q)$. Niech
$\lambda_s:\Delta^{p+q}\to\Delta^p\times\Delta^q$
będzie mapą afiniczną na te wierzchołki. Oznaczmy przez
$\varepsilon_s$ znak permutacji przenoszącej uporządkowane
$p$ kroki poziome przed $q$ pionowymi do słowa $s$.
Definiujemy
$\mathrm{sh}(\sigma\otimes\tau)
=\sum_s\varepsilon_s(\sigma\times\tau)\lambda_s$.
Wszystkie wewnętrzne ściany triangulacji produktu
$\Delta^p\times\Delta^q$ pojawiają się dwukrotnie ze znakami
przeciwnymi: odpowiadają przeplotom różniącym się pojedynczym
zamienieniem sąsiednich kroków. Zewnętrzne ściany dają
$\mathrm{sh}(\partial\sigma\otimes\tau)$ i
$(-1)^p\mathrm{sh}(\sigma\otimes\partial\tau)$.
To wyjaśnia zarówno wystąpienie znaku, jak i tożsamość
$\partial\mathrm{sh}=\mathrm{sh}\partial$.

\emph{Krok 3: lemat o rozszerzaniu homotopii na modelach.}
Załóżmy, że $F,G:C\to D$ są mapami łańcuchowymi i $C_n$
jest wolna na wskazanych generatorach-modelach. Mamy skonstruować
$K_n:C_n\to D_{n+1}$ tak, aby $\partial K+K\partial=F-G$.
Przy konstrukcji indukcyjnej na generatorze $u$ stopnia $n$
wyraz $z=(F-G)u-K(\partial u)$ jest cyklem, ponieważ
$\partial z=(F-G)(\partial u)-(F-G)(\partial u)+K\partial^2u=0$.
Jeśli modelowa przestrzeń wartości ma zerową homologię
zredukowaną w stopniu $n$, można wybrać $K(u)$ z
$\partial K(u)=z$. Przy $n=0$ wymagamy dodatkowo zerowej
augmentacji $z$. Budując $K$ w modelu, a następnie
przenosząc go przez mapy $\sigma_\#$, otrzymujemy naturalność.

Uzasadnijmy wspomniane znikanie: dla wypukłego $V$ wybieramy
$v\in V$ i stożkujemy sympleksy ku $v$. Operator $S$ spełnia
$\partial S+S\partial=\id-\eta\epsilon$, gdzie
$\epsilon:C_0(V)\to\Z$ jest augmentacją i
$\eta:\Z\to C_0(V)$ wysyła $1$ na stały sympleks w $v$.
W dodatnich stopniach każdy cykl jest więc brzegiem.
Analogicznie tensor dwóch takich kompleksów ma kontrakcję:
na elementach jednorodnych działa
$S_X\otimes\id+(\eta\epsilon)_X\otimes S_Y$,
ze standardową konwencją znaków dla map stopnia $1$.
Po podstawieniu otrzymujemy
$\partial S_{\otimes}+S_{\otimes}\partial
=\id-(\eta\epsilon)_X\otimes(\eta\epsilon)_Y$.

\emph{Krok 4: $\mathrm{sh}\,\mathrm{AW}\simeq\id$.}
Dla uniwersalnego sympleksu w iloczynie wybieramy
$\Delta^n\xrightarrow{\Delta_n}\Delta^n\times\Delta^n$,
$t\mapsto(t,t)$. Jego obrazy w dowolnym $X\times Y$
powstają przez mapy $\sigma_X\times\sigma_Y$.
W stopniu zero obie mapy są równe. Jeśli $K$ znamy na
ścianach $\Delta_n$, to
$z=(\mathrm{sh}\,\mathrm{AW}-\id)(\Delta_n)-K(\partial\Delta_n)$
jest cyklem w $C_n(\Delta^n\times\Delta^n)$. Model jest
wypukły, więc cykl ma wypełnienie $K(\Delta_n)$.
Dla $n=0$ augmentacja zniknęła już na starcie.
Przenosząc wybrane wypełnienia przez
$(\sigma_X\times\sigma_Y)_\#$, otrzymujemy żądaną
homotopię naturalną. Istotne jest, że pracujemy z
kompleksem \emph{produktu} $\Delta^n\times\Delta^n$.

\emph{Krok 5: $\mathrm{AW}\,\mathrm{sh}\simeq\id$.}
Tym razem modele dla generatorów $\sigma\otimes\tau$ to
$\id_{\Delta^p}\otimes\id_{\Delta^q}$, a grupy docelowe
modelowej indukcji to
$C_*(\Delta^p)\otimes C_*(\Delta^q)$.
Dla $p+q=0$ mapy są identyczne. W wyższym stopniu
$z=(\mathrm{AW}\,\mathrm{sh}-\id)
(\id_{\Delta^p}\otimes\id_{\Delta^q})-K\partial
(\id_{\Delta^p}\otimes\id_{\Delta^q})$ jest cyklem.
Kontrakcja tensorowa z kroku 3 wypełnia ten cykl.
Przeniesienie przez $\sigma_\#\otimes\tau_\#$ daje $K$
na wszystkich generatorach. Indukcja dowodzi tożsamości
$\partial K+K\partial=\mathrm{AW}\,\mathrm{sh}-\id$.
\end{proof}

\begin{twierdzenie}[Künneth dla homologii singularnej o współczynnikach całkowitych]
\label{thm:kunneth-v4}
Niech $X,Y$ będą dowolnymi przestrzeniami topologicznymi, a $n\geq0$.
Istnieje naturalny względem obu przestrzeni krótki ciąg dokładny
\[
\begin{gathered}
0\longrightarrow
\bigoplus_{p+q=n}H_p(X;\Z)\otimes_\Z H_q(Y;\Z)
\xrightarrow{\ \times\ }H_n(X\times Y;\Z)\\
\xrightarrow{\ \kappa\ }
\bigoplus_{p+q=n-1}\operatorname{Tor}^{\Z}_1(H_p(X;\Z),H_q(Y;\Z))
\longrightarrow0.
\end{gathered}
\]
Ciąg rozszczepia się jako ciąg grup abelowych, ale na ogół nie istnieje
rozszczepienie naturalne. Nie zakładamy, że którakolwiek grupa homologii
jest skończenie generowana.
\end{twierdzenie}
\begin{proof}[Część algebraiczna i standardowa naturalność ciągu]
Dowód podzielimy na cztery kroki, oddzielając istnienie ciągu
\emph{naturalnego} od konstrukcji jego \emph{nienaturalnego} rozszczepienia.

\emph{1. Rozkład kompleksu wolnego.}
Niech $C$ będzie nieujemnie stopniowanym kompleksem wolnych grup
abelowych. Piszemy $Z_k=\ker(\partial_k:C_k\to C_{k-1})$ oraz
$B_k=\operatorname{im}(\partial_{k+1}:C_{k+1}\to C_k)$.
Każda podgrupa grupy wolnej abelowej jest wolna, również
przy nieskończonej randze. Wobec tego $B_{k-1}$ jest grupą
projektywną i ciąg
$0\to Z_k\to C_k\xrightarrow{\partial_k}B_{k-1}\to0$
rozszczepia się. Wybierzmy rozszczepienia
$s_{k-1}:B_{k-1}\to C_k$. Otrzymujemy izomorfizm
$C_k\cong Z_k\oplus B_{k-1}$, przy którym
$\partial_k(z,b)=(b,0)\in Z_{k-1}\oplus B_{k-2}$.
Stąd $C$ jest (niekanonicznie) izomorficzny z sumą prostą
kompleksów $P_k=(B_k\hookrightarrow Z_k)$ w stopniach
$k+1,k$. Każdy $P_k$ jest dwustopniową wolną rezolwentą
$H_k(C)=Z_k/B_k$. Ten argument nie korzysta z twierdzenia
strukturalnego o skończenie generowanych grupach abelowych.

\emph{2. Tensor z dowolnym kompleksem wolnym.}
Niech $D$ będzie drugim kompleksem grup wolnych.
Dla ustalonego $k$ użyjmy dokładnego ciągu kompleksów
$0\to Z_k\otimes D\to P_k\otimes D\to (B_k\otimes D)[1]\to0$,
gdzie $[1]$ oznacza przesunięcie stopnia homologicznego o jeden.
Dokładność wynika z tego, że w obu stopniach $P_k$
rozkład rozpatruje bezpośrednio kompleks stożka
$B_k\otimes D\to Z_k\otimes D$.
Ponieważ $Z_k,B_k$ są wolne, mamy
$H_j(Z_k\otimes D)\cong Z_k\otimes H_j(D)$ i analogicznie
dla $B_k$. Długi ciąg homologii stożka daje zatem
\[
\begin{split}
0\to&H_k(C)\otimes H_j(D)\longrightarrow
H_j(P_k\otimes D)\\
\longrightarrow&\operatorname{Tor}^{\Z}_1(H_k(C),H_{j-1}(D))\to0.
\end{split}
\]
W powyższym zapisie $P_k$ na chwilę traktujemy bez przesunięcia,
jako kompleks w stopniach $1,0$; przy wstawieniu go w stopnie
$k+1,k$ oba indeksy homologiczne należy zwiększyć o $k$.
Identyfikacja jądra z $\operatorname{Tor}_1$ pochodzi z wolnej
rezolwenty $0\to B_k\to Z_k\to H_k(C)\to0$.
W tym kroku nie wybraliśmy żadnego rozkładu samych grup homologii.

\emph{3. Naturalność ciągu.}
Aby uniknąć pozornej naturalności wybranych rozszczepień
$C_k\cong Z_k\oplus B_{k-1}$, definiujemy lewą strzałkę
bezpośrednio jako iloczyn zewnętrzny cykli
$[c]\otimes[d]\mapsto[c\otimes d]$.
Opis powyższy pokazuje, że jest ona monomorfizmem,
a jej kokernel to wskazana suma grup $\operatorname{Tor}_1$.
Naturalna strzałka $\kappa$ wymaga osobnej konstrukcji:
można ją otrzymać z homomorfizmów łączących w ciągach
dokładnych filtracji tensora wolnych kompleksów, równoważnie
z naturalnej ciągu Künnetha. Same niekanoniczne rozkłady
$C_k=Z_k\oplus B_{k-1}$ nie dowodzą jej naturalności.
W tym miejscu przyjmujemy standardowy lemat o naturalności
konstrukcji $\kappa$; pozostałe kroki wykazują dokładność,
identyfikację członów i nienaturalne rozszczepienie.

\emph{4. Rozszczepienie bez skończonej generacji.}
Przeprowadźmy rozkład z kroku 1 \emph{dla obu} kompleksów.
Tensor rozdziela się nad dowolnymi sumami prostymi:
$C\otimes D\cong\bigoplus_{k,l}P_k(C)\otimes P_l(D)$
z uwzględnieniem przesunięć.
Dla dwóch wolnych rezolwent
$P(A):F_1\hookrightarrow F_0$ i
$P(B):G_1\hookrightarrow G_0$ kompleks tensorowy ma
zaledwie trzy stopnie $0,1,2$. Jego homologia w stopniu $2$
zanika: różniczka
$F_1\otimes G_1\to(F_0\otimes G_1)\oplus(F_1\otimes G_0)$
jest injektywna, bo $F_1\hookrightarrow F_0$ i $G_1$
jest płaska. Homologia w stopniu $0$ to $A\otimes B$,
a w stopniu $1$ to $\operatorname{Tor}_1(A,B)$:
aby to sprawdzić, rozpatrzmy dwustopniową filtrację
bikompleksu $P(A)\otimes P(B)$ według stopnia $P(A)$.
Dla ustalonego wolnego $F_i$ kompleks pionowy
$F_i\otimes(G_1\hookrightarrow G_0)$ ma homologię
$F_i\otimes B$ w stopniu $0$ i zero w stopniu $1$,
bo $F_i$ jest płaska. Zatem homologia całego tensoru
jest homologią dwustopniowego kompleksu
$F_1\otimes B\to F_0\otimes B$, a jego grupy
w stopniach $1,0$ są z definicji odpowiednio
$\operatorname{Tor}_1(A,B)$ i $A\otimes B$. Zatem homologia \emph{każdego}
$P_k(C)\otimes P_l(D)$ ma odrębne składowe
$H_k(C)\otimes H_l(D)$ w stopniu $k+l$ i
$\operatorname{Tor}_1(H_k(C),H_l(D))$ w stopniu $k+l+1$.
Ponieważ homologia komutuje z sumami prostymi kompleksów,
wybrane rozkłady obu kompleksów dają poszukiwane
rozszczepienie, choć zależne od wyborów.

Na koniec stosujemy twierdzenie Eilenberga--Zilbera:
$C_*(X\times Y)$ jest naturalnie homotopijnie równoważny
$C_*(X)\otimes C_*(Y)$; grupy łańcuchów singularnych są wolne
bez żadnych założeń o skończoności. Wobec tego powyższy
ciąg algebraiczny daje tezę, a jego lewa strzałka jest
klasycznym iloczynem zewnętrznym.
\end{proof}
\begin{przyklad}[Nieskończenie generowana grupa i brak problemu ze skończonością]
Dla $A=\bigoplus_{j\ge1}\Z/2^j$ wybierzmy dwustopniową wolną
rezolwentę
$\bigoplus_{j\ge1}\Z\xrightarrow{\oplus_j\,2^j}
 \bigoplus_{j\ge1}\Z$.
Wówczas dla $B=\Z/2$ mamy
$A\otimes B\cong\bigoplus_{j\ge1}\Z/2$ oraz
$\operatorname{Tor}^{\Z}_1(A,B)\cong\bigoplus_{j\ge1}\Z/2$.
Wyprowadzenie wynika bezpośrednio z tensora wolnych rezolwent;
nie wolno użyć tu rozkładu dla \emph{skończenie} generowanych
grup w zastępstwie dowodu ogólnego.
\end{przyklad}


\begin{przyklad}[Künneth: torsja w dwóch sąsiednich wymiarach]
\label{prz:kunneth-rp2-razy-rp2}
Niech $P=\mathbb{RP}^2$ i $X=P\times P$.
Z Przykładu~\ref{prz:homologia-rp2} wiemy, że
$H_0(P;\Z)=\Z$, $H_1(P;\Z)=\Z/2\Z$,
a wszystkie grupy $H_k(P;\Z)$ dla $k\geq2$ są zerowe.
W stopniu $2$ lewy wyraz ciągu Künnetha to
\[
 \bigoplus_{p+q=2}H_p(P)\otimes H_q(P)
 \cong(\Z/2\Z)\otimes(\Z/2\Z)
 \cong\Z/2\Z.
\]
Prawy wyraz jest zerowy: sumy
$\operatorname{Tor}_1(H_p(P),H_q(P))$ dla $p+q=1$
mają czynnik wolny $\Z$, więc znikają.
Stąd $H_2(X;\Z)\cong\Z/2\Z$.

W stopniu $3$ lewy wyraz jest zerowy,
ponieważ trzeba by użyć $H_2(P)$ albo $H_3(P)$.
Z prawej strony jedyny niezerowy składnik to
\[
 \operatorname{Tor}_1^\Z(H_1(P),H_1(P))
 =\operatorname{Tor}_1^\Z(\Z/2\Z,\Z/2\Z)
 \cong\Z/2\Z.
\]
Otrzymujemy więc $H_3(X;\Z)\cong\Z/2\Z$.
W innych dodatnich wymiarach homologia znika
z wyjątkiem $H_1(X;\Z)\cong(\Z/2\Z)^2$.
Zwróćmy uwagę na mechanizm: klasa w stopniu $2$
pochodzi z iloczynu tensorowego, ale klasa
w stopniu $3$ pochodzi \emph{wyłącznie z funktora
$\operatorname{Tor}_1$}.

Z twierdzenia o współczynnikach uniwersalnych
dla kohomologii otrzymujemy dodatkowo
\[
 H^3(X;\Z)\cong
 \operatorname{Ext}^1_\Z(H_2(X),\Z)
 \cong\Z/2\Z,
 \qquad H^4(X;\Z)\cong\Z/2\Z.
\]
Jest to zgodne z kohomologiczną wersją twierdzenia
Künnetha: w stopniu $3$ pojawia się składnik
$\operatorname{Tor}_1(H^2(P),H^2(P))$, natomiast
w stopniu $4$ iloczyn tensorowy tych klas.
\end{przyklad}

\begin{twierdzenie}[Künneth dla kohomologii skończonych CW-kompleksów]
Niech $X,Y$ będą skończonymi CW-kompleksami. Dla $n\ge0$
istnieje naturalny krótki ciąg dokładny
\[
0\to\bigoplus_{p+q=n}H^p(X;\Z)\otimes H^q(Y;\Z)
\xrightarrow{\times}H^n(X\times Y;\Z)
\to\bigoplus_{p+q=n+1}\operatorname{Tor}^{\Z}_1
(H^p(X;\Z),H^q(Y;\Z))\to0.
\]
Rozszczepienie jest na ogół nienaturalne.
\end{twierdzenie}
\begin{proof}
Skończony rozkład CW daje kompleksy komórkowe z wolnych
grup skończonej rangi. Komórki produktu mają postać
$e^p\times e^q$, a wzór na ich granice to
$\partial(e^p\times e^q)=\partial e^p\times e^q
+(-1)^p e^p\times\partial e^q$.
Dla skończonych wolnych grup kanoniczna mapa
$\operatorname{Hom}(A,\Z)\otimes\operatorname{Hom}(B,\Z)
\to\operatorname{Hom}(A\otimes B,\Z)$ jest izomorfizmem.
Otrzymujemy więc iloczyn tensorowy komórkowych kompleksów
\emph{kołańcuchowych}, a do niego stosujemy powyższe
twierdzenie algebraiczne, zmieniając kierunek stopniowania.
Stąd po prawej występuje warunek $p+q=n+1$.
Lewa strzałka jest iloczynem zewnętrznym
$\alpha\times\beta=\mathrm{pr}_X^*\alpha\smile\mathrm{pr}_Y^*\beta$;
identyfikację potwierdza sprawdzenie na bazach dualnych komórek.
\end{proof}

\begin{przyklad}
Dla $\R P^2$ mamy $H^0(\R P^2;\Z)=\Z$,
$H^2(\R P^2;\Z)=\Z/2$ i brak innych dodatnich grup
kohomologii. Wzór kohomologiczny w stopniu $3$ daje
$H^3(\R P^2\times\R P^2;\Z)\cong\Z/2$ wyłącznie dzięki
$\operatorname{Tor}_1(\Z/2,\Z/2)$.
\end{przyklad}

\begin{uwaga}[Stopniowa przemienność iloczynu kubkowego]
Wcześniejszy dowód stopniowej przemienności można domknąć
konstrukcją z kroku 3. Niech $D=\mathrm{AW}\circ\Delta_\#$,
gdzie $\Delta:X\to X\times X$ jest diagonalą, a
$\tau(a\otimes b)=(-1)^{pq}b\otimes a$ dla stopni $p,q$.
Mapy $D$ i $\tau D$ są naturalnymi mapami łańcuchowymi
$C_*(X)\to C_*(X)\otimes C_*(X)$, zgodnymi w stopniu zero.
Ich różnicę w stopniu $n$ na modelu $\Delta^n$ wypełniamy
indukcyjnie kontrakcją tensoru
$C_*(\Delta^n)\otimes C_*(\Delta^n)$, identycznie jak w
dowodzie Eilenberga--Zilbera. Powstaje homotopia
$D-\tau D=\partial K+K\partial$. Po zastosowaniu
$\alpha\otimes\beta$ do tego wzoru, dla kocykli stopni
$p,q$, dostajemy
$\alpha\smile\beta-(-1)^{pq}\beta\smile\alpha=\delta\gamma$.
To podaje brakujący argument bez fałszywego twierdzenia o
przemienności na poziomie samych kołańcuchów.
\end{uwaga}


\section{Związek z kohomologią de Rhama}

\begin{definicja}
Niech $M$ będzie gładką rozmaitością, a $\Omega^k(M)$ przestrzenią
gładkich form różniczkowych stopnia $k$. Różniczka zewnętrzna
$\mathrm d:\Omega^k(M)\to\Omega^{k+1}(M)$ spełnia $\mathrm d^2=0$.
Kohomologią de Rhama nazywamy
$H^k_{\mathrm{dR}}(M)=\ker(\mathrm d:\Omega^k(M)\to\Omega^{k+1}(M))
/\operatorname{im}(\mathrm d:\Omega^{k-1}(M)\to\Omega^k(M))$,
przyjmując $\Omega^{-1}(M)=0$. Pełne konstrukcje i przykłady
pojawią się w następnym rozdziale.
\end{definicja}

\begin{twierdzenie}[Twierdzenie de Rhama, postać porównawcza]
Niech $M$ będzie gładką rozmaitością, a $n\geq0$. Całkowanie
gładkich form $n$-stopnia po gładkich sympleksach singularnych
indukuje naturalny izomorfizm
$H^n_{\mathrm{dR}}(M)\cong H^n(M;\R)$.
Izomorfizmy te zachowują iloczyn zewnętrzny po lewej i iloczyn
kubkowy po prawej na poziomie pierścieni kohomologii.
\end{twierdzenie}

\begin{przyklad}[Generator kohomologii de Rhama okręgu]
\label{prz:de-rham-okrag-explicit}
Niech $S^1=\{(x,y)\in\R^2:x^2+y^2=1\}$.
Ograniczenie do $S^1$ formy
\[
 \omega=\frac{x\,\mathrm dy-y\,\mathrm dx}{2\pi}
\]
jest zamknięte, ponieważ każda forma stopnia $2$
na jednowymiarowej rozmaitości znika.
Dla dodatnio zorientowanej parametryzacji
$\gamma(t)=(\cos(2\pi t),\sin(2\pi t))$, $t\in[0,1]$,
mamy $\gamma^*\omega=\mathrm dt$, a zatem
\[
 \int_{S^1}\omega=\int_{[0,1]}\gamma^*\omega=1.
\]
Forma $\omega$ nie jest dokładna:
gdyby $\omega=\mathrm df$ na $S^1$,
to z twierdzenia Stokesa całka po zamkniętym
okręgu wynosiłaby $0$. Jej klasa jest
generatorem $H^1_{\mathrm{dR}}(S^1)\cong\R$
w sensie normalizacji całki.
Po porównaniu de Rhama odpowiada ona
singularnemu kocyklowi, którego wartość
na klasie fundamentalnej okręgu wynosi $1$.
\end{przyklad}

\begin{uwaga}
Dla formy zamkniętej $\omega$ jej klasie odpowiada kocykl
$\sigma\mapsto\int_{\Delta^n}\sigma^*\omega$.
Twierdzenie Stokesa gwarantuje, że jest to kocykl,
a forma dokładna daje kobrzeg. Pełne uzasadnienie, że
konstrukcja jest izomorfizmem, korzysta z lematu Poincarégo
i ciągu Mayera--Vietorisa; jest rozwinięte w następnym
rozdziale. Zgodność iloczynów na kohomologii wynika z
porównania przekątnej singularnej z iloczynem zewnętrznym,
nie z punktowej równości wartości na pojedynczym sympleksie.
\end{uwaga}

